% GENERATED FILE. Edit canonical lessons, then run npm run generate:books.
% canonical-source-hash: fnv1a64:0037ac7e90bfb96f
% canonical-lessons: ZH-C281-biru, ZH-W281-ru, ZH-C281-yinwei-suoyi, ZH-W281-suo, ZH-C281-yinci, ZH-C281-lizi, ZH-W281-one-reason, ZH-C281-practice

\chapter{For Example}
\label{ch:zh-for-example}
\hlchaptermodality{281}

\begin{chapteropening}
\textbf{By the end of this chapter:} \emph{I can give a reason and an example.}

Giving a reason with \zh{因为} … \zh{所以} … in one sentence, for instance, for this reason and an example, and writing one line of your own with a reason in it, built on two new written characters: \zh{如} and \zh{所}.
\end{chapteropening}

\section[\texorpdfstring{\zh{比如} (bǐrú)}{比如 (bǐrú)}]{\zh{比如} (bǐrú) — for instance}

\label{lesson:ZH-C281-biru}



\begin{quote}
A new chapter, and it opens on the word that brings in an example.

\begin{itemize}
\raggedright
  \item \emph{Recall:} say \zh{我觉得这个意见很好} — \emph{wǒ juéde zhège yìjiàn hěn hǎo} — \textbf{R1}, the chapter you have finished
  \item \emph{Recall:} say \zh{这么} — \emph{zhème}, so — \textbf{R1}
  \item \emph{Recall:} say \zh{比} — \emph{bǐ}, than — \textbf{R3}, durable
\end{itemize}
\end{quote}

\subsection*{You'll want to know: \zh{比如}}
\textbf{\zh{比如}} — \emph{bǐrú} — \textquotedblleft{}for instance; such as\textquotedblright{}.

\begin{quote}
\textbf{\zh{我喜欢中国菜} \zh{比如包子}} \emph{wǒ xǐhuan Zhōngguó cài, bǐrú bāozi} — I like Chinese food, steamed buns for instance.
\end{quote}

\begin{quote}
\textbf{\zh{我喜欢喝茶} \zh{比如中国茶}} \emph{wǒ xǐhuan hē chá, bǐrú Zhōngguó chá} — I like tea, Chinese tea for instance.
\end{quote}

The next lesson puts the second character in your hand. The first is the \zh{比} you write.

\begin{cousinweb}
\textbf{\zh{比如}} (\emph{bǐrú}) is built from \zh{比} 'to compare' and \zh{如} 'like', both inherited from Middle Chinese. The Guangyun lists \zh{比} under five readings, of which \emph{pjijX} is glossed 'to compare', and \zh{如} under two, \emph{nyo} and \emph{nyoH}; it glosses \emph{nyo} with several words, among them 'to resemble' (Guangyun, in Baxter's transcription). 'For instance' for the pair is the word's own sense.
\end{cousinweb}

\begin{grammarlens}[title={What you've built}]
\zh{比如} comes after the general point and before the example: \zh{中国菜}, then \zh{比如包子}.

It is the \zh{比} of comparing, \zh{比} as in \zh{我比他大}, with an example set beside the point.
\end{grammarlens}

\subsection*{Your turn}
\begin{itemize}
\raggedright
  \item \emph{Say it:} \emph{bǐrú}
  \item \emph{Say it:} \emph{bǐrú}, once more
  \item \emph{Say it:} \zh{我喜欢中国菜} \zh{比如包子} — \emph{wǒ xǐhuan Zhōngguó cài, bǐrú bāozi} — I like Chinese food, steamed buns for instance.
  \item \emph{Say it:} \zh{我喜欢喝茶} \zh{比如中国茶} — \emph{wǒ xǐhuan hē chá, bǐrú Zhōngguó chá} — I like tea, Chinese tea for instance.
  \item \emph{Recall:} say \zh{比}, then \zh{比如}
\end{itemize}

\begin{tcolorbox}[breakable,colback=teal!4,colframe=teal!35!black,title={Before you move on}]
What does \zh{比如} mean? (\textquotedblleft{}For instance; such as\textquotedblright{}.) Where does it go? (\textbf{After the general point, before the example}.)
\end{tcolorbox}
\section[\texorpdfstring{\zh{如} (rú)}{如 (rú)}]{\zh{如} — \zh{女}, then \zh{口}}

\label{lesson:ZH-W281-ru}



\begin{quote}
The word first, then the character inside it.

\begin{itemize}
\raggedright
  \item \emph{Recall:} say \emph{bǐrú} — \textbf{R1}, one lesson back, before it decays
\end{itemize}

\begin{quote}
\zh{好} \zh{妈} \zh{口}
\end{quote}

\emph{Read it:} each one aloud, with nothing said first — \emph{hǎo}, \emph{mā}, \emph{kǒu}
\end{quote}

\subsection*{Script — \zh{女}, then \zh{口}}
\hlblockfigure{\detokenize{figures/ZH-W281-ru-filmstrip.pdf}}{How \zh{如} is written, stroke by stroke}

\begin{quote}
\zh{如}
\end{quote}

\textbf{Two pieces, and you write both already.}

\textbf{\zh{女}, three strokes}, narrow, on the left, as in \zh{好} and \zh{妈}: down to the left and turning down to the right without lifting; a left-falling stroke from the upper right; and a horizontal that rises as it crosses.

\textbf{\zh{口}, three strokes}, on the right: the left side, the top that turns down the right side, the bottom bar.

3 + 3 = six.

\subsection*{Your turn}
\begin{itemize}
\raggedright
  \item \emph{Write it:} \zh{如}, all six, with the model in front of you
  \item \emph{Cover:} write it again from memory, and count the pieces, not the strokes
\end{itemize}

\begin{tcolorbox}[breakable,colback=teal!4,colframe=teal!35!black,title={Before you move on}]
How many strokes in \zh{如}? (\textbf{Six}.) Which piece comes first? (\textbf{\zh{女}, on the left}.) Where else does \zh{女} stand on the left? (\textbf{In \zh{好} and \zh{妈}}.)

Source: \href{https://raw.githubusercontent.com/chanind/hanzi-writer-data/68d10a4b21150cae5e1ebbd223eed289cf32d90c/data/\%E5\%A6\%82.json}{Hanzi Writer Data \zh{如}, pinned stroke paths and medians}.
\end{tcolorbox}
\section[\texorpdfstring{\zh{因为}…\zh{所以}… (yīnwèi … suǒyǐ …)}{因为…所以… (yīnwèi … suǒyǐ …)}]{\zh{因为}…\zh{所以}… — because, so}

\label{lesson:ZH-C281-yinwei-suoyi}



\begin{quote}
Two words you have had for a long time, joined into one frame.

\begin{itemize}
\raggedright
  \item \emph{Read it:} \zh{比如}, aloud — \emph{bǐrú}
  \item \emph{Recall:} say \zh{因为} — \emph{yīnwèi}, because — \textbf{R3}, durable
  \item \emph{Recall:} say \emph{suǒyǐ}, so — \textbf{R3}
  \item \emph{Recall:} say \zh{为什么} — \emph{wèishénme}, why — \textbf{R3}
\end{itemize}
\end{quote}

\begin{grammarlens}[title={Grammar Lens: \zh{因为} … \zh{所以} …, a reason and what follows}]
\begin{quote}
\textbf{\zh{因为} opens the reason, \zh{所以} opens what follows from it, and together they make one sentence.}
\end{quote}

\begin{quote}
\textbf{\zh{因为下雨了} \zh{所以我们没去}} \emph{yīnwèi xiàyǔ le, suǒyǐ wǒmen méi qù} — Because it rained, we didn't go. \textbf{\zh{因为他病了} \zh{所以没来上课}} \emph{yīnwèi tā bìng le, suǒyǐ méi lái shàngkè} — He was ill, so he didn't come to class.
\end{quote}

You have had \zh{因为} and \emph{suǒyǐ} as separate words; now they frame one sentence, and the next lesson puts \zh{所} in your hand so the frame prints whole.

English keeps one of the two words; Mandarin is happy with both.

To answer \zh{为什么}, \zh{因为} alone is enough: \textbf{\zh{为什么没去} — \zh{因为下雨了}} \emph{wèishénme méi qù — yīnwèi xiàyǔ le}.
\end{grammarlens}

\subsection*{Your turn}
Say these aloud:

\begin{itemize}
\raggedright
  \item \zh{因为下雨了} \zh{所以我们没去} — \emph{yīnwèi xiàyǔ le, suǒyǐ wǒmen méi qù} — Because it rained, we didn't go.
  \item \zh{因为他病了} \zh{所以没来上课} — \emph{yīnwèi tā bìng le, suǒyǐ méi lái shàngkè} — He was ill, so he didn't come to class.
  \item \zh{为什么没去} — \zh{因为下雨了} — \emph{wèishénme méi qù — yīnwèi xiàyǔ le} — Why didn't you go? Because it rained.
  \item why you are tired today, with \zh{因为} … \zh{所以} …
\end{itemize}

\begin{tcolorbox}[breakable,colback=teal!4,colframe=teal!35!black,title={Before you move on}]
Which word opens the reason? (\textbf{\zh{因为}}.) Which opens what follows? (\textbf{\zh{所以}}.) How do you answer \zh{为什么} with a reason alone? (\textbf{\zh{因为}, and the reason}.)
\end{tcolorbox}
\section[\texorpdfstring{\zh{所} (suǒ)}{所 (suǒ)}]{\zh{所} — two four-stroke sides}

\label{lesson:ZH-W281-suo}



\begin{quote}
The sentence first, then the character it was waiting for.

\begin{itemize}
\raggedright
  \item \emph{Recall:} say \zh{因为下雨了} \zh{所以我们没去} — \emph{yīnwèi xiàyǔ le, suǒyǐ wǒmen méi qù} — \textbf{R1}, one lesson back, before it decays
\end{itemize}

\begin{quote}
\zh{房} \zh{以}
\end{quote}

\emph{Read it:} each one aloud, with nothing said first — \emph{fáng}, \emph{yǐ}
\end{quote}

\subsection*{Script — two four-stroke sides}
\hlblockfigure{\detokenize{figures/ZH-W281-suo-filmstrip.pdf}}{How \zh{所} is written, stroke by stroke}

\begin{quote}
\zh{所}
\end{quote}

\textbf{Two sides of four strokes each, and both are new in this shape.}

\textbf{The left side, four strokes}: a short, flat left-falling stroke at the top; a long left-falling stroke down the left side; a horizontal that turns down at its right end; and a horizontal under it. It looks like the top of \zh{房}, but \zh{房} opens with a dot, and this opens with a short left-falling stroke.

\textbf{The right side, four strokes}: a flat left-falling stroke at the top; a long left-falling stroke that drops down and then curves to the lower left; a horizontal across it; and a long vertical down to the foot.

4 + 4 = eight.

\subsection*{Your turn}
\begin{itemize}
\raggedright
  \item \emph{Write it:} \zh{所}, all eight, with the model in front of you
  \item \emph{Cover:} write it again from memory, and count the pieces, not the strokes
\end{itemize}

\begin{tcolorbox}[breakable,colback=teal!4,colframe=teal!35!black,title={Before you move on}]
How many strokes in \zh{所}? (\textbf{Eight}.) How does the left side begin? (\textbf{With a short left-falling stroke, not the dot of \zh{房}}.) Which stroke comes last? (\textbf{The long vertical on the right}.)

Source: \href{https://raw.githubusercontent.com/chanind/hanzi-writer-data/68d10a4b21150cae5e1ebbd223eed289cf32d90c/data/\%E6\%89\%80.json}{Hanzi Writer Data \zh{所}, pinned stroke paths and medians}.
\end{tcolorbox}
\section[yīncǐ]{yīncǐ — for this reason}

\label{lesson:ZH-C281-yinci}



\begin{quote}
Before the new one: the character you wrote last, and the word it sits in.

\begin{itemize}
\raggedright
  \item \emph{Read it:} \zh{所以}, aloud — \emph{suǒyǐ}
  \item \emph{Recall:} say \zh{因为下雨了} \zh{所以我们没去} — \emph{yīnwèi xiàyǔ le, suǒyǐ wǒmen méi qù} — \textbf{R1}, two lessons back
\end{itemize}
\end{quote}

\subsection*{You'll want to know: yīncǐ}
\emph{yīncǐ} — \textquotedblleft{}for this reason; and so\textquotedblright{}.

\begin{quote}
\emph{tiānqì bù hǎo, yīncǐ wǒmen méi qù} — The weather was bad, and so we didn't go.
\end{quote}

Its second character comes later; for now it is a word for the ear.

\begin{cousinweb}
\emph{yīncǐ} is built from \zh{因} and cǐ 'this', both inherited from Middle Chinese, \emph{'jin} and \emph{tshjeX} (Guangyun, in Baxter's transcription). The dictionary glosses the first 'to rely on' and 'a cause', and the second 'to stop'; 'this' is not given there, and 'for this reason' for the pair is the word's own sense.
\end{cousinweb}

\begin{grammarlens}[title={What you've built}]
\emph{suǒyǐ} is the everyday so. \emph{yīncǐ} belongs to writing and careful speech, and it usually stands on its own, with no \zh{因为} in front.

It opens the result, after the reason: \emph{tiānqì bù hǎo, yīncǐ wǒmen méi qù}.
\end{grammarlens}

\subsection*{Your turn}
\begin{itemize}
\raggedright
  \item \emph{Say it:} \emph{yīncǐ}
  \item \emph{Say it:} \emph{yīncǐ}, once more
  \item \emph{Say it:} \emph{tiānqì bù hǎo, yīncǐ wǒmen méi qù} — The weather was bad, and so we didn't go.
  \item \emph{Recall:} say \emph{suǒyǐ}, then \emph{yīncǐ}: the everyday so, then the careful one
\end{itemize}

\begin{tcolorbox}[breakable,colback=teal!4,colframe=teal!35!black,title={Before you move on}]
What does \emph{yīncǐ} mean? (\textquotedblleft{}For this reason\textquotedblright{}.) Where does it belong? (\textbf{In writing and careful speech}.)
\end{tcolorbox}
\section[lìzi]{lìzi — an example}

\label{lesson:ZH-C281-lizi}



\begin{quote}
Before the new one: last lesson's word, and the word that brings an example in.

\begin{itemize}
\raggedright
  \item \emph{Recall:} say \emph{yīncǐ}, for this reason — \textbf{R1}, one lesson back
  \item \emph{Read it:} \zh{比如}, aloud — \emph{bǐrú}
\end{itemize}
\end{quote}

\subsection*{You'll want to know: lìzi}
\emph{lìzi} — \textquotedblleft{}an example\textquotedblright{}.

\begin{quote}
\emph{qǐng shuō yí ge lìzi} — Please give an example.
\end{quote}

Its first character comes later; for now it is a word for the ear.

\begin{cousinweb}
\emph{lìzi} is built from lì 'a precedent' and the ending \zh{子}, both inherited from Middle Chinese, \emph{ljejH} and \emph{tsiX} (Guangyun, in Baxter's transcription). The dictionary glosses the first 'to compare' and 'all'; 'an example' is the modern word's own sense.
\end{cousinweb}

\begin{grammarlens}[title={What you've built}]
\zh{比如} brings an example in; \emph{lìzi} is the example itself, a noun you can count: \emph{yí ge lìzi}.

So a teacher asks \emph{qǐng shuō yí ge lìzi}, and you answer with \zh{比如} ….
\end{grammarlens}

\subsection*{Your turn}
\begin{itemize}
\raggedright
  \item \emph{Say it:} \emph{lìzi}
  \item \emph{Say it:} \emph{lìzi}, once more
  \item \emph{Say it:} \emph{qǐng shuō yí ge lìzi} — Please give an example.
  \item \emph{Say it:} \zh{比如包子} — \emph{bǐrú bāozi} — Steamed buns, for instance.
  \item \emph{Recall:} say \zh{比如}, then \emph{lìzi}: bringing one in, then the example itself
\end{itemize}

\begin{tcolorbox}[breakable,colback=teal!4,colframe=teal!35!black,title={Before you move on}]
What does \emph{lìzi} mean? (\textquotedblleft{}An example\textquotedblright{}.) How do you ask for one? (\emph{qǐng shuō yí ge lìzi}.)
\end{tcolorbox}
\section[(one reason)]{One reason — a line of your own with \zh{因为} and \zh{所以}}

\label{lesson:ZH-W281-one-reason}



\begin{quote}
Someone asks you a question you have surely been asked before: \zh{你为什么学汉语}?

This lesson answers it, in one line of your own.

\begin{itemize}
\raggedright
  \item \emph{Recall:} say \emph{lìzi}, an example
\end{itemize}
\end{quote}

\subsection*{Writing — controlled composition}
Write \textbf{one line} that answers \zh{你为什么学汉语}. Use the frame: \zh{因为} and your reason, then \zh{所以} and \zh{我学汉语}. If you want to add an example, \zh{比如} brings it in at the end.

A word whose characters you have not written goes in pinyin, inside the line, with its tone marks: \zh{中国菜} is all characters you write, but \emph{kěyǐ} has one you do not write yet, so it goes in pinyin. Every character comes from the ones you have learned, each in the order you learned it, each in the same size of square.

No Chinese commas or full stops: leave a space between the two halves, as this book prints them.

Write the line. Then read it back: is \zh{因为} in front of the reason, and \zh{所以} in front of what follows?

\begin{tcolorbox}[breakable,colback=teal!4,colframe=teal!35!black,title={Before you move on}]
Which half of your line could stand alone as the answer to \zh{为什么}? (\textbf{The half with \zh{因为}}.)

Which character did you write for the first time this chapter, inside your line? (\textbf{\zh{所}, in \zh{所以}}.)
\end{tcolorbox}
\section[Practice]{\zh{因为我喜欢中国} \zh{所以学汉语}}

\label{lesson:ZH-C281-practice}



\begin{quote}
The chapter closes here, with the question your own line answered.

\begin{itemize}
\raggedright
  \item \emph{Recall:} say \emph{lìzi}, an example — \textbf{R2}
  \item \emph{Recall:} say \emph{yīncǐ}, for this reason — \textbf{R2}
  \item \emph{Read it:} \zh{如}, aloud — \emph{rú}, the second half of \zh{比如}
\end{itemize}
\end{quote}

\subsection*{The exchange}
\begin{quote}
\textbf{A:} \zh{你为什么学汉语}
\end{quote}

\begin{quote}
\emph{nǐ wèishénme xué Hànyǔ} — Why are you learning Chinese?
\end{quote}

\begin{quote}
\textbf{B:} \zh{因为我喜欢中国} \zh{所以学汉语}
\end{quote}

\begin{quote}
\emph{yīnwèi wǒ xǐhuan Zhōngguó, suǒyǐ xué Hànyǔ} — I like China, so I'm learning Chinese.
\end{quote}

\begin{quote}
\textbf{A:} \zh{你喜欢中国的什么}
\end{quote}

\begin{quote}
\emph{nǐ xǐhuan Zhōngguó de shénme} — What do you like about China?
\end{quote}

\begin{quote}
\textbf{B:} \zh{中国菜} \zh{中国茶}
\end{quote}

\begin{quote}
\emph{Zhōngguó cài, Zhōngguó chá} — Chinese food, Chinese tea.
\end{quote}

\begin{quote}
\textbf{A:} \emph{nǐ néng shuō yí ge lìzi ma} — Can you give an example?
\end{quote}

\begin{quote}
\textbf{B:} \zh{比如包子} \zh{我很喜欢}
\end{quote}

\begin{quote}
\emph{bǐrú bāozi, wǒ hěn xǐhuan} — Steamed buns, for one. I love them.
\end{quote}

\begin{quote}
\textbf{A:} \emph{wǒ yě shì, yīncǐ wǒ xiǎng qù Zhōngguó} — Me too, and so I want to go to China.
\end{quote}

\zh{因为} … \zh{所以} … gives the reason and the result in one breath; \zh{比如} then backs it up. The line with \emph{yīncǐ} sounds a little careful, as that word does.

\begin{tcolorbox}[breakable,skin=enhanced,colback=orange!6,colframe=orange!45!black,boxrule=0.5pt,arc=1mm,left=6pt,right=6pt,top=4pt,bottom=4pt,fonttitle=\bfseries,title={Read this}]
\emph{Read it:} the exchange through once without stopping, then answer the question below
\end{tcolorbox}

\begin{grammarlens}[title={What you've built — a reason and an example}]
A reason, what follows from it, and something to show it:

\begin{itemize}
\raggedright
  \item \textbf{\zh{因为下雨了} \zh{所以我们没去}} \emph{yīnwèi xiàyǔ le, suǒyǐ wǒmen méi qù} — the reason and the result.
  \item \emph{tiānqì bù hǎo, yīncǐ wǒmen méi qù} — the careful version.
  \item \textbf{\zh{中国菜} \zh{比如包子}} \emph{Zhōngguó cài, bǐrú bāozi} — an example brought in.
  \item \emph{qǐng shuō yí ge lìzi} — asking for one.
\end{itemize}

\textbf{Listen.} \emph{Nǐ wèishénme xué Hànyǔ?}

\textbf{Speak.} Give two reasons for something you do every week, each with \zh{因为} … \zh{所以} ….

\emph{Write it:} \zh{所以} and \zh{比如}
\end{grammarlens}

\subsection*{Your turn}
\begin{itemize}
\raggedright
  \item \emph{Say it:} \zh{因为我病了} \zh{所以我今天不去上课} — \emph{yīnwèi wǒ bìng le, suǒyǐ wǒ jīntiān bú qù shàngkè} — I'm ill, so I'm not going to class today.
  \item \emph{Say it:} \zh{为什么} — \zh{因为我喜欢中国} — \emph{wèishénme — yīnwèi wǒ xǐhuan Zhōngguó} — Why? Because I like China.
  \item \emph{Say it:} the whole exchange, both parts
  \item \emph{Write it:} \zh{因为我病了} \zh{所以没去}
\end{itemize}

\begin{tcolorbox}[breakable,colback=teal!4,colframe=teal!35!black,title={Before you move on}]
Say because it rained, we didn't go. (\textbf{\zh{因为下雨了} \zh{所以我们没去}}.) Say steamed buns, for instance. (\textbf{\zh{比如包子}}.) Ask for an example. (\emph{qǐng shuō yí ge lìzi}.)
\end{tcolorbox}
